\documentclass[a4paper]{arrowhead}

\usepackage[yyyymmdd]{datetime}
\usepackage{etoolbox}
\usepackage[utf8]{inputenc}
\usepackage[T1]{fontenc}
\usepackage{multirow}
\usepackage{longtable}
\usepackage{tikz}
\usepackage{adjustbox}
\AtBeginEnvironment{longtable}{\small}
\BeforeBeginEnvironment{tikzpicture}{\begin{adjustbox}{max width=\textwidth}}
\AfterEndEnvironment{tikzpicture}{\end{adjustbox}}
\usetikzlibrary{positioning,shapes.multipart,arrows.meta}

\renewcommand{\dateseparator}{-}
\setlength{\emergencystretch}{3em}
\lstset{language=HTTP,breakatwhitespace=false,columns=fullflexible}

%% Special references
\newcommand{\mref}[2][]{{\textcolor{ArrowheadPurple}{\hyperref[sec:model:#2]{\ifstrempty{#1}{#2}{#1}}}}}
\newcommand{\oref}[1]{{\textcolor{ArrowheadBlue}{\hyperref[sec:operations:#1]{#1}}}}
\newcommand{\pdef}[1]{{\textcolor{ArrowheadGrey}{#1\label{sec:model:primitives:#1}\label{sec:model:primitives:#1s}\label{sec:model:primitives:#1es}}}}
\newcommand{\pref}[1]{{\textcolor{ArrowheadGrey}{\hyperref[sec:model:primitives:#1]{#1}}}}

% as template_sd.tex; the input and output arguments are passed already formatted (\mref or \pref)
\newrobustcmd\osubsection[3]{
  \addtocounter{subsection}{1}
  \addcontentsline{toc}{subsection}{\protect\numberline{\thesubsection}operation \textcolor{ArrowheadBlue}{#1}}
  \subsection*{
    \thesubsection\quad
    operation
    \textcolor{ArrowheadBlue}{#1}
    (\notblank{#2}{#2}{})
    \notblank{#3}{: #3}{}
  }
  \label{sec:operations:#1}
}
\newrobustcmd\msubsection[2]{
  \addtocounter{subsection}{1}
  \addcontentsline{toc}{subsection}{\protect\numberline{\thesubsection}#1 \textcolor{ArrowheadPurple}{#2}}
  \subsection*{\thesubsection\quad#1 \textcolor{ArrowheadPurple}{#2}}
  \label{sec:model:#2} \label{sec:model:#2es}
}
%%

\usepackage{xurl}
\begin{document}
\sloppy

%% Arrowhead Document Properties
\ArrowheadTitle{identity}
\ArrowheadServiceID{identity}
\ArrowheadType{Service Description}
\ArrowheadTypeShort{SD}
\ArrowheadVersion{5.3.0}
\ArrowheadDate{2026-11-09}
\ArrowheadAuthor{Jerker Delsing}
\ArrowheadStatus{REVIEW DRAFT}
\ArrowheadContact{jerker.delsing@ltu.se}
\ArrowheadFooter{\href{www.arrowhead.eu}{www.arrowhead.eu}}
\ArrowheadSetup
%%

%% Front Page
\begin{center}
  \vspace*{1cm}
  \huge{\arrowtitle}

  \vspace*{0.2cm}
  \LARGE{\arrowtype}
  \vspace*{1cm}

  \Large{Service ID: \textit{"\arrowid"}}

  \vspace*{\fill}

  % Front Page Image

  \vspace*{1cm}
  \vspace*{\fill}

  % Front Page Abstract
  \begin{abstract}
    This document describes the \texttt{identity} service, version 1.0.0, produced by the Authentication core system of Arrowhead (documentation edition 5.3.0).
  \end{abstract}

  \vspace*{1cm}
\end{center}

\newpage
%%

%% Table of Contents
\tableofcontents
\newpage
%%

\section{Documentation edition and implementation scope}
This is the 5.3.0 documentation edition, dated 2026-11-09. It carries forward the system and service contracts from the retained 5.2.2 documentation and identifies the current primary implementation sources. The former \path|documentation/v5.2.2| folder is retained unchanged.

Document version 5.3.0 is distinct from executable, dependency, registered system/service and wire-contract versions. Values still defined by the source (including 5.2.2 artefact versions, 1.0.0 service versions and \path|SysMonitor-AMQP-IDD-5.2.2| payload identifiers) are preserved. This documentation update does not change runtime behaviour or establish a new implementation release. Historical verification statements describe the tests recorded in the inherited baseline; this revision validates document compilation and references.

\section{Overview}
\label{sec:overview}

This document describes the \texttt{identity} service, which provides authentication of systems: a system logs in with its credentials and obtains an identity token, which it presents to other systems; providers verify such tokens.
It is produced by the Authentication system \cite{auth-sysd} in service version 1.0.0.

The rest of this document is organized as follows.
In Section \ref{sec:operations}, we describe the abstract message operations provided by the service.
In Section \ref{sec:model}, we end the document by presenting the data types used by the mentioned operations.
The concrete interface is specified in the HTTP/TLS/JSON IDD \cite{auth-identity-http-idd}. The MQTT binding is specified in the MQTT/TLS/JSON IDD \cite{auth-identity-mqtt-idd}.

\subsection{Significant Prior Art}

Arrowhead 4.x authenticated systems only by X.509 certificates. The 5.x identity service adds credential-based login and identity tokens, used where certificates are not available or as a second factor.

\subsection{How This Service Is Meant to Be Used}

A system calls \texttt{identity-\allowbreak{}login} and sends the returned token as \texttt{Authorization: Bearer IDENTITY-\allowbreak{}TOKEN/\allowbreak{}/\allowbreak{}\textless{}token\textgreater{}}. Providers with the \texttt{outsourced} authentication policy call \texttt{identity-\allowbreak{}verify} to learn which system a token belongs to and whether it is a sysop. \texttt{identity-\allowbreak{}logout} ends the session, \texttt{identity-\allowbreak{}change-\allowbreak{}credentials} replaces the password.

\subsection{Important Delimitations and Dependencies}
\label{sec:delimitations}

Authentication keeps one session per system: a new login replaces the previous token of that system. In this workspace login, logout and credential change over HTTPS are accepted only with the client certificate of the same system, and every token is bound to the certificate it is used with (token-to-certificate binding). The service does not authorize anything.

\newpage

\section{Service Interface}
\label{sec:operations}

This section describes the interfaces to the \texttt{identity} service.
In particular, each subsection names an abstract operation, an input type and an output type, in that order.
The input type is named inside parentheses, while the output type is preceded by a colon.
Input and output types are only denoted when accepted or returned, respectively, by the interface in question.

All abstract data types named in this section are defined in Section \ref{sec:model}.

\begin{figure}[ht!]
  \centering
  \begin{tikzpicture}[font=\small]
    \node[draw=ArrowheadBlue, very thick, rectangle split, rectangle split parts=2, align=center,
          rectangle split part fill={ArrowheadBlue!12,white}, text width=9cm, inner sep=6pt] (b)
      {\guillemotleft{}Service\guillemotright{}\\\textbf{identity}
        \nodepart{second} \texttt{identity-\allowbreak{}login}\\ \texttt{identity-\allowbreak{}logout}\\ \texttt{identity-\allowbreak{}change-\allowbreak{}credentials}\\ \texttt{identity-\allowbreak{}verify}};
  \end{tikzpicture}
  \caption{SysML-style block of the \texttt{identity} service and its operations.}
  \label{fig:SD}
\end{figure}

The following interface operations are available.

\osubsection{identity-login}{\mref{IdentityRequest}}{\mref{IdentityLoginResponse}}

Checks the credentials and opens a session. Returns the identity token and its expiry (\texttt{identity.\allowbreak{}token.\allowbreak{}duration}, 300 s in the default configuration). A previous token of the system becomes invalid.

\osubsection{identity-logout}{\mref{IdentityRequest}}{}

Checks the credentials and closes the session of the system.

\osubsection{identity-change-credentials}{\mref{IdentityChangeRequest}}{}

Checks the current credentials and replaces them with the new ones.

\osubsection{identity-verify}{\pref{Token}}{\mref{IdentityVerifyResponse}}

Checks an identity token. Returns whether it is valid and, if so, the system, its sysop flag, login and expiry time. The caller must itself present a valid identity token (only \texttt{login}, \texttt{logout} and \texttt{change} are token-free).

\newpage

\section{Information Model}
\label{sec:model}

Here, all data objects that can be part of the \texttt{identity} service are listed in alphabetic order.
Note that each subsection, which describes one type of object, begins with the \textit{struct} keyword, which is used to denote a collection of named fields, each with its own data type.
As a complement to the explicitly defined types in this section, there is also a list of implicit primitive types in Section \ref{sec:model:primitives}, which are used to represent things like hashes and identifiers.

\msubsection{struct}{IdentityChangeRequest}

Changes the credentials of the requesting system.


\noindent\begin{tabularx}{\textwidth}{| >{\raggedright\arraybackslash}p{3.9cm} | >{\raggedright\arraybackslash\small}p{5.3cm} | X |} \hline
\rowcolor{gray!33} Field & Type & Description \\ \hline
system\allowbreak{}Name & \pref{Name} & Name of the system, PascalCase (e.g. \texttt{TemperatureSensor}). \\ \hline
credentials & \pref{Map}$<$\allowbreak{}\pref{String}, \pref{String}$>$ & Credentials of the authentication method, for \texttt{PASSWORD} the key \texttt{password}. \\ \hline
new\allowbreak{}Credentials & \pref{Map}$<$\allowbreak{}\pref{String}, \pref{String}$>$ & The new credentials. \\ \hline
\end{tabularx}



\msubsection{struct}{IdentityLoginResponse}

The identity token issued at login.


\noindent\begin{tabularx}{\textwidth}{| >{\raggedright\arraybackslash}p{3.9cm} | >{\raggedright\arraybackslash\small}p{5.3cm} | X |} \hline
\rowcolor{gray!33} Field & Type & Description \\ \hline
token & \pref{Token} & The identity token, used as \texttt{Bearer IDENTITY-\allowbreak{}TOKEN/\allowbreak{}/\allowbreak{}\textless{}token\textgreater{}}. \\ \hline
expiration\allowbreak{}Time & \pref{DateTime} & UTC time when the identity token expires. \\ \hline
\end{tabularx}



\msubsection{struct}{IdentityRequest}

System name and credentials.


\noindent\begin{tabularx}{\textwidth}{| >{\raggedright\arraybackslash}p{3.9cm} | >{\raggedright\arraybackslash\small}p{5.3cm} | X |} \hline
\rowcolor{gray!33} Field & Type & Description \\ \hline
system\allowbreak{}Name & \pref{Name} & Name of the system, PascalCase (e.g. \texttt{TemperatureSensor}). \\ \hline
credentials & \pref{Map}$<$\allowbreak{}\pref{String}, \pref{String}$>$ & Credentials of the authentication method, for \texttt{PASSWORD} the key \texttt{password}. \\ \hline
\end{tabularx}



\msubsection{struct}{IdentityVerifyResponse}

Result of verifying an identity token.


\noindent\begin{tabularx}{\textwidth}{| >{\raggedright\arraybackslash}p{3.9cm} | >{\raggedright\arraybackslash\small}p{5.3cm} | X |} \hline
\rowcolor{gray!33} Field & Type & Description \\ \hline
verified & \pref{Boolean} & True when the token is valid. \\ \hline
system\allowbreak{}Name & \pref{Name} & System the token belongs to. \\ \hline
sysop & \pref{Boolean} & Whether that system is a system operator. \\ \hline
login\allowbreak{}Time & \pref{DateTime} & UTC time of the login. \\ \hline
expiration\allowbreak{}Time & \pref{DateTime} & UTC time when the identity token expires. \\ \hline
\end{tabularx}



\subsection{Primitives}
\label{sec:model:primitives}

Types and structures mentioned throughout this document that are assumed to be available to implementations of this service.
The concrete interpretations of each of these types and structures must be provided by any IDD document claiming to implement this service.

\noindent\begin{tabularx}{\textwidth}{| p{4.6cm} | X |} \hline
\rowcolor{gray!33} Type & Description \\ \hline
\pdef{Boolean} & One out of \texttt{true} or \texttt{false}. \\ \hline
\pdef{DateTime} & Pinpoints a specific moment in time (UTC). \\ \hline
\pdef{Map} & A set of key-value pairs. \\ \hline
\pdef{Name} & A string identifier that is intended to be both human and machine-readable. \\ \hline
\pdef{String} & A chain of characters. \\ \hline
\pdef{Token} & An opaque string used as credential. \\ \hline
\end{tabularx}

\newpage

\bibliographystyle{IEEEtran}
\bibliography{bibliography}

\newpage

\section{Revision History}
\subsection{Amendments}

\noindent\begin{tabularx}{\textwidth}{| p{1cm} | p{3cm} | p{2cm} | X | p{4cm} |} \hline
\rowcolor{gray!33} No. & Date & Version & Subject of Amendments & Author \\ \hline

1 & 2026-10-22 & 5.2.2 & Document created from the Arrowhead LaTeX templates for the 5.2.2 Java Spring implementation, including the local extensions of this workspace. & Jerker Delsing \\ \hline
2 & 2026-11-09 & 5.3.0 & New documentation edition; primary-source references and Go/Rust implementation scope; 5.2.2 retained. & Documentation maintenance \\ \hline
\end{tabularx}

\subsection{Quality Assurance}

\noindent\begin{tabularx}{\textwidth}{| p{1cm} | p{3cm} | p{2cm} | X |} \hline
\rowcolor{gray!33} No. & Date & Version & Approved by \\ \hline

1 & & \arrowversion  &  \\ \hline

\end{tabularx}

\end{document}
